\documentclass{article}
\usepackage{graphicx} % Required for inserting images
\usepackage[ruled, vlined, linesnumbered]{algorithm2e}
\usepackage{amsmath}
\usepackage{parskip}
\usepackage{cancel}

\title{32 bit Modular Arithmetic and SIMD Vectorisation}
\author{Duc Vinh Nguyen, Henry Zheng}
\date{March 2026}
\begin{document}
\maketitle

\section{Introduction}
This project explores and optimize modular reductions for prime numbers  $p < 2^{32}$ using SIMD Vectorisation (AVX2, AVX512, NEON) along with various different techniques and algorithmes.

\textbf{Objective: } The aim of this study is to be able to realize an optimized implementation with minimal performance cost for modular reductions.

\section{Fundamental Cost of Modular Reduction}
The most naive way to compute $a \pmod{n}$ is to perform the operation :

\begin{equation}
    a - \left\lfloor \frac{a}{n} \right\rfloor \cdot n
\end{equation}

However, on modern processor architectures (Intel, AMD,\dots), integer division is non-pipelined and has significantly higher computational cost than other operations like additions and multiplications.

Since equation (1) requires a divsion and substraction operation, its considered highly inefficient from both a software and hardware perspective.

\section{Shoup's Modular Multiplication}
Shoup's algorithm is centraled around precomputation as a work around the costly division operations

\begin{algorithm}[H]
    \caption{Shoup's modular multiplication algorithm}
    
    % Input and Output definitions
    \SetKwInOut{Input}{Input}
    \SetKwInOut{Output}{Output}

    % The content
    \Input{$A, B \in [0, p[$, we suppose that $p < \beta$ \\
           \hspace{1.2cm} precompute $B' = \lfloor B\beta/p \rfloor$}
    \Output{$R = AB \bmod p$}
    
    \BlankLine
    
    % Steps
    $Q \leftarrow \lfloor AB'/\beta \rfloor$ \quad ;\\
    $R \leftarrow (AB - Qp) \bmod \beta$ \quad ;
    
    \If{$R \ge p$}{
        $R \leftarrow R - p$\;
    }
\end{algorithm}

\begin{description}
    \item[Specifications:]\hfill
    \begin{itemize}
        \item $\beta$: For this case is $2^{32}$
        \item A: An integer over 32 bits ( $<\beta = 2^{32}$)
        \item B: An integer over 32 bits ( $<\beta = 2^{32}$)
        \item p: A prime number over 32 bits ( $< \beta = 2^{32}$)
        \item Q: The quotient
        \item R: The remainder
    \end{itemize}
\end{description}

\subsection{Procomputational Logic}
The precomputed value $B'$ is defined over 32 bits as:
\[B'= \lfloor (B \cdot 2^{32}) / p \rfloor \]

When computing the $a \cdot b \pmod{p}$ :
\begin{equation*}
  A \cdot B \pmod{p} = A \cdot B -  \left\lfloor \frac{A \cdot B}{p} \right\rfloor \cdot p
\end{equation*}
We can rewrite this: 
\begin{align*}
  A \cdot B \pmod{p} &= A \cdot B -  \left\lfloor \frac{A \cdot \left\lfloor (\frac{B \cdot 2^{32}}{p}) \right\rfloor}{2^{32}} \right\rfloor \cdot p \\
  &\approx A \cdot B -  \left\lfloor \frac{A \cdot B'}{2^{32}} \right\rfloor \cdot p \\
\end{align*}

Using the precomputed value B' to approxiamtely calculate $\left\lfloor \frac{A \cdot B}{p} \right\rfloor$ improves on the naive method by 2 ways:
\begin{itemize}
  \item $B'$ is computed with 1 division and can be stored to be used in other modular reduction provided the same B
  \item $\frac{A \cdot B'}{2^{32}}$ has only 1 multiplication operation and a bit shift to the right of 32 bits. On a CPU, both multiplication and bit shift operations can be done efficiently compared to division
\end{itemize}

\subsection{Fundamental Principal}
Given x the quotient of 2 32 bits integers: \\
The proof of this algorithm based on the non decimal value of ($\lfloor x \rfloor$): 
$$x - 1 < \lfloor x \rfloor \le x$$
We can rewrite as:
$$0 \le x - \lfloor x \rfloor < 1$$

\subsection{Basic Reduction B = 1:}

\subsection{Demonstration:}
The algorithm uses the precomputation $B' = \lfloor B\beta/p \rfloor$ such that:
\[0 \le \frac{B\beta}{p} - \left\lfloor \frac{B\beta}{p} \right\rfloor < 1 \]
\begin{equation}
    0 \le \frac{B\beta}{p} - B' < 1
\end{equation}

Using the same logic, the algorithm estimates the quotient $Q = \lfloor AB'/\beta \rfloor$  such that:
\[0 \le \frac{AB'}{\beta} - \left\lfloor\frac{AB'}{\beta}\right\rfloor < 1\]
\begin{equation}
    0 \le \frac{AB'}{\beta} - Q < 1
\end{equation}

We now prouve that remainder $R = (A \cdot B -Q \cdot p)$ is within the interval [0, 2p[:

We firstly multiply the equation (1) by $Ap/\beta$:
\[0 \le \left( \frac{B\beta}{p} - B' \right) \cdot \frac{Ap}{\beta} < 1 \cdot \frac{Ap}{\beta}\]
\begin{equation}
    0 \le AB - \frac{B'Tp}{\beta} < \frac{Bp}{\beta}
\end{equation}

Similarly, we multiply the equation (2) by $p$:
\[0 \le \left( \frac{AB'}{\beta} - Q \right) \cdot p < 1 \cdot p\]
\begin{equation}
    0 \le \frac{AB'p}{\beta} - Qp < p
\end{equation}

By adding the equation (3) and (4), we deduce:
\[0 \le (AB - \cancel{\frac{AB'p}{\beta}}) + (\cancel{\frac{AB'p}{\beta}} - Qp) < \frac{Ap}{\beta} + p\]
Simplifying this result, we are left with:
\begin{align*}
    0 \le AB - Qp &< p + \frac{Ap}{\beta} \\
    0 \le AB - Qp &< p \cdot (1 + \frac{A}{\beta}) 
\end{align*}

By constraint, $A \in [0, p[$ and $p < \beta$ thus $A < \beta$. Therefore we can rewrite:
\begin{align*}
    0 \le AB - Qp &< p \cdot (1 + \frac{\beta}{\beta}) \\
    0 \le AB - Qp &< p \cdot (1 + 1)\\
    0 \le AB - Qp &< 2p \\
    0 \le R &< 2p
\end{align*}

We have now prouven that $A \in [0, 2p[$ the final if statement of the helps correct the remainder. If $R > p$ the 
modular reduction needs 1 more step $R - p$, if not $R$ is the correct result


\subsection{SIMD Vectorization}
We decided to create 2 versions for implementating a high performance Shoup algorithm
\subsubsection{Intel \_mm256\_mul\_epu32}
\subsubsection{Intel \_mm256\_mullo\_epi32}
\subsection{Performance}
\subsubsection{Choice of unrolling}
\subsubsection{Simplification when B = 1}
\subsubsection{Overall Performance}
\subsection{Avantages:}

\end{document}